\begin{exercise}{126}
Show that the category of marked simplicial sets is canonically enriched in
simplicial sets.

\begin{proof}
For marked simplicial sets $A=(X,S)$ and $B=(Y,T)$, take the hom-object
$\mathrm{Map}^{\sharp}(A,B)$ of Remark~3.2.15:
\[
 \mathrm{Map}^{\sharp}(A,B)_n
 =\operatorname{Hom}_{\mathrm{sSet}_+}
   ((\Delta^n)^{\sharp}\times A,B).
\]
Thus its simplices are maps $F\colon\Delta^n\times X\to Y$ satisfying
$F_1(e,s)\in T$ whenever $s\in S$, with simplicial operators given by
precomposition in the first factor.
For $C=(Z,U)$, define composition in degree $n$ by
\[
 (G\circ F)(t,x)=G(t,F(t,x)).
\]
If $s\in S$, then $F_1(e,s)\in T$, hence
$G_1(e,F_1(e,s))\in U$; therefore this composite preserves markings.
The formula commutes with precomposition in $\Delta^n$, so it defines
a map of simplicial sets
\[
 \mathrm{Map}^{\sharp}(B,C)\times\mathrm{Map}^{\sharp}(A,B)
 \longrightarrow\mathrm{Map}^{\sharp}(A,C).
\]
The identity of $A$ defines the unit vertex, whose degeneracy in degree $n$
is the marked projection $\Delta^n\times X\to X$.
Associativity and both unit laws follow strictly from these formulas.
In degree zero the hom-sets and composition are exactly those of
$\mathrm{sSet}_+$.
Everything is defined by products, projections and composition, naturally
in marked maps, so this enrichment is canonical.
\end{proof}
\end{exercise}
